\chapter{第十四讲\quad 刑罚的体系}
\label{lecture:14}
\noindent{\small\hyperref[ref:lecture:14]{查阅本讲扩展规则}}\par

\section{第一节\quad 主刑}
\paragraph{刑种、期限与执行}
第33条规定的主刑为管制、拘役、有期徒刑、无期徒刑、死刑。以 \(s\) 表示刑种，\(L\) 表示宣告刑期，\(D\) 表示判决执行前依法应予折抵的先行羁押日数。单个刑期的一般约束及折抵是
\[
\begin{array}{c|c|c}
s & L\text{ 的范围} & \text{折抵日数}\\ \hline
\text{管制} & [3\text{个月},\,2\text{年}] & 2D\\
\text{拘役} & [1\text{个月},\,6\text{个月}] & D\\
\text{有期徒刑} & [6\text{个月},\,15\text{年}] & D
\end{array}
\]
分别依据第38至47条。有期徒刑期限的第50、69条例外另算。刑期从判决执行日起算；月份、年份的期限计算按法律口径，不能通用固定 \(30\) 日换算。管制实行社区矫正，可以有禁止令；拘役由公安机关就近执行；有期、无期徒刑在监狱或其他执行场所执行。无期徒刑没有宣告的固定月数，减刑或假释另有条件。

\paragraph{例题推导}
甲依法被判管制六个月，先行羁押十日。先识别 \(s=\text{管制}\)；依第41条选择系数 \(2\)；得 \(2\times10=20\) 日可折抵。六个月仍是宣告期限，折抵的是应执行部分。若改判拘役六个月，则依据第44条，系数改为 \(1\)，十日只能折十日。若试图判一罪管制两年六个月，超过第38条上限；两罪并罚后的执行刑期另按第69条，上限三年。

\paragraph{死刑的排除条件与死缓分支}
令 \(J\) 表示罪行极其严重且法条允许适用死刑，\(U\) 表示犯罪时未满十八周岁，\(P\) 表示审判时怀孕，\(O\) 表示审判时已满七十五周岁，\(C\) 表示以特别残忍手段致人死亡。第48、49条给出必要条件
\[
 \text{死刑候选}\ \Longrightarrow\
 J\land\neg U\land\neg P\land(\neg O\lor C).
\]
条件满足后，仍须依案件情节裁量。应判死刑且无须立即执行时，可以宣告死刑缓期二年执行。依第50、51条，二年期间从判决确定之日起算。无故意犯罪的通常出口是期满减为无期徒刑；确有重大立功的法定出口为期满减为二十五年有期徒刑。故意犯罪且情节恶劣的，报最高人民法院核准后执行死刑。故意犯罪而未执行死刑的，死缓期间重新计算并报最高人民法院备案。死缓减为有期徒刑的刑期从死缓期满之日起算。对死缓累犯及因第50条第二款列举犯罪被判死缓者，法院依法可以同时限制减刑。死刑立即执行原则上由最高人民法院核准，依法由最高人民法院判决的除外；死缓可以由高级人民法院判决或核准。

反例：甲作案时十七岁，审判时十九岁。即使 \(J\) 已满足，\(U\) 为真使必要条件失败，立即执行与死缓均被排除。对审判时七十六岁者，只有第49条特别例外可能移除年龄排除，移除排除也不等于应判死刑。

\section{第二节\quad 附加刑}
\paragraph{按附加后果分别建立约束}
第34条列罚金、剥夺政治权利、没收财产，可以独立适用；第35条允许对犯罪的外国人独立或附加适用驱逐出境。第52、53条的罚金以犯罪情节决定数额，有缴纳期限及强制缴纳、随时追缴情形；遇不能抗拒的灾祸等确有困难，须经法院裁定，方可能延期、酌减或免除。第59、60条的没收财产以本人所有部分或全部为对象，没收全部时须保留本人及其扶养家属必需的生活费用；家属所有或应有财产不得没收。没收财产以前的正当债务需要从被没收财产偿还且债权人提出请求时，应当偿还。第36条还规定赔偿被害人的民事责任在指定冲突情形下优先。

剥夺政治权利所涉权利依第54条分别核验；第55至58条的常见期限分支为
\[
\begin{array}{l|l}
\text{情形} & \text{期限及执行}\\ \hline
\text{一般有限期} & [1,5]\text{年}\\
\text{管制附加} & \text{与管制期限相同并同时执行}\\
\text{判处死刑、无期徒刑} & \text{终身}\\
\text{死缓或无期徒刑依法减为有期徒刑} & [3,10]\text{年}
\end{array}
\]
危害国家安全犯罪依法应当附加；故意杀人、强奸、放火等严重破坏社会秩序犯罪可以附加。附加于徒刑、拘役时，附加刑的刑期通常从主刑执行完毕之日或假释日起算，其效力在主刑执行期间也适用。

\paragraph{例题推导}
甲被判管制八个月并附加剥夺政治权利。输入为“管制附加”而非一般分支；由第55条第二款，附加刑亦为八个月且同步执行，所以不受通常的一年下限限制。若改为有期徒刑附加，八个月便低于通常下限；还须重新确定起算点。若另判没收财产，不能把配偶拥有的份额并入 \(M=\{a:a\text{ 属于被告本人}\}\)。

\section{第三节\quad 非刑罚处罚措施}
\paragraph{对象集合与适用入口}
第37条可在犯罪情节轻微、不需要判刑时免予刑事处罚，并依案件情况适用训诫、赔礼道歉等措施。第37条之一的从业禁止须先检验“被判处刑罚”以及利用职业便利犯罪或违反职业特定义务犯罪，再由法院结合预防再犯罪等需要决定；相关职业的限制从刑罚执行完毕日或假释日起算，范围为三至五年。其他法律、行政法规另有职业禁限规定的从其规定。区分刑中禁止令与刑后从业禁止，看对象、起点及条文。

第64条按照财物属性选择处置：
\[
\begin{array}{l|l}
\text{经证明的属性} & \text{处置路径}\\ \hline
\text{违法所得} & \text{追缴或责令退赔}\\
\text{被害人的合法财产} & \text{及时返还}\\
\text{违禁品} & \text{没收}\\
\text{供犯罪所用且属犯罪人本人财物} & \text{没收}
\end{array}
\]
最后一行的两项事实是合取，证明其一不足以适用该行；另有法律依据的处置另论。

\paragraph{例题推导}
甲用朋友不知情借出的普通相机实施犯罪，偷来的首饰仍在甲手中，甲还持有依法认定的违禁品。首饰属于被害人的合法财产，依第64条追回后返还；违禁品依独立分支没收。相机虽是犯罪工具，但不是甲本人财物，所以不能仅依“供犯罪所用的本人财物”这一分支没收。若相机依法证实属于甲，合取两项同时为真，结论随之改变。

\section*{记忆与闭卷复述}
记“主刑看种限，附加看对象，措施看财物”。先写主刑三组通常期限与折抵系数 \(2,1,1\)；再从作案时、审判时、判决确定后三条时间线复述死刑门槛；最后把“个人合法财产、违法所得、被害人合法财产、违禁品、本人犯罪工具”分成五张财物卡逐一判定。每次记住一条通常规则后，配写一条例外。
\begin{enumerate}
 \item \textbf{自测。}审判时已满七十五周岁、犯罪时已成年且未使用特别残忍手段致人死亡，可否判死缓？\quad\textbf{核对。}不适用死刑，死缓也在排除范围。
 \item \textbf{自测。}甲利用职业便利犯罪被判刑，法院决定三年职业禁止，甲其后假释。三年从哪天算？\quad\textbf{核对。}依第37条之一从假释日起算；缓刑、假释的考验期限是另一套期限。
\end{enumerate}
